% Source-aligned English for printed p.163.
\SourcePage{163}

\LayerTarget{tohoku:claim:3.4.2:H1-bijection:surjectivity}
if $L$ is locally isomorphic to $\overline F$ (for one readily sees that a
sheaf locally isomorphic to a flabby sheaf is flabby). It is an epimorphism.
Indeed, let $(h_{ij})$ be a 1-cocycle of $H$ relative to an open covering
$(U_i)$ of $X$. Since $H^1(X,\overline H)$ consists only of the neutral
element, there exist $\overline h_i\in\Gamma(U_i,\overline H)$ such that
$h_{ij}=\overline h_i^{-1}\overline h_j$. On the other hand, the
$\overline h_i$ can plainly be lifted to sections $\overline g_i$ of
$\overline G$ (it suffices to recall the definitions of $\overline G$ and
$\overline H$!). Put $p_{ij}=\overline g_i^{-1}\overline g_j$. The section
of $H$ over $U_{ij}$ defined by $p_{ij}$ is then $h_{ij}$; hence
$p_{ij}\in\Gamma(U_{ij},P)$. Thus $(p_{ij})$ is a 1-cocycle of $P$ that
defines the given 1-cocycle upon passage to the quotient.

\LayerTarget{tohoku:construction:3.4.2:associated-P-bundle}
It follows from the preceding result that the fiber space $E$ is isomorphic
to the fiber space associated with a fiber space $Q$ having structure sheaf
$P$; moreover, $Q$ is determined up to isomorphism.

\LayerTarget{tohoku:action:3.4.2:P-actions}
Let $P=\overline F G\subset\overline G$ act on $\overline G$ by inner
automorphisms. Then $F$ remains stable under these actions, since
$\overline F$ acts trivially on $\overline F$ and $F$ is invariant in $G$.
Thus $P$ acts naturally on $\overline F$. Moreover, $F$ is stable under the
actions of $P$, and the resulting actions of $P$ on $F$ are precisely those
obtained by composing $P\longrightarrow H$ with the natural representation
of $H$ by actions on $F$. It follows that the associated sheaf $F^E$ also
identifies with the associated sheaf $F^Q$. The actions of $P$ also descend
to the quotient $F'=\overline F/F$; the representations of $P$ by germs of
automorphisms of $F,\overline F,F'$ will all be denoted by $\sigma$. We now
have an exact sequence of homomorphisms of sheaves:

\LayerTarget{tohoku:eq:3.4.2:crossed-extension}
\[
e\longrightarrow G\longrightarrow P\mathop{\longrightarrow}^{u}F'
\longrightarrow e.
\]

\LayerTarget{tohoku:map:3.4.2:u}
The homomorphism $G\longrightarrow P=\overline F G$ is the inclusion and
hence a homomorphism of sheaves of groups. The map
$u:P\longrightarrow F'$ sends a product $\overline f g$, with
$\overline f\in\overline F(x)$ and $g\in G(x)$, to the class of $\overline f$
in $F'(x)$. This is well-defined because $G\cap\overline F=F$.

\LayerTarget{tohoku:claim:3.4.2:u-properties}
This map \textit{does not preserve} the multiplicative structures in general,
but it satisfies the following conditions: (i) $u$ is surjective, and two
elements of $P$ have the same image if and only if they are congruent under
the right action of $G$, that is, if they define the same element of $P/G$;
(ii) $u$ is a crossed homomorphism from $P$ to $F'$, where $P$ acts on $F'$
as indicated above. Thus, in every fiber $P(x)$,
$u(e)=e$ and $u(pp')=u(p)\sigma(p)u(p')$.

\LayerTarget{tohoku:construction:3.4.2:dQ-obstruction}
From this situation and from the fiber space $Q$ with structure sheaf $P$ we
will derive an element $d(Q)\in H^1(X,{F'}^Q)$, where ${F'}^Q$ is the sheaf
associated with $Q$ and the given actions of $P$ on $F'$, in such a way that
$d(Q)=0$ if and only if the class $c(Q)\in H^1(X,P)$ of $Q$ lies in the image
of $H^1(X,G)$. Proposition 3.4.2 will then be proved, modulo the immediate
verifications of naturality.

\LayerTarget{tohoku:reduction:3.4.2:E-via-Q}
Indeed, the resulting condition is necessary and sufficient for the class
$c(E)\in H^1(X,H)$ to lie in the image of $H^1(X,G)$, as is immediately seen
from the commutative diagram

\TohokuContinuation{164}
